\documentclass[12pt]{article}
\usepackage{amsmath,amssymb,geometry,enumitem,booktabs}
\geometry{margin=1in}
\setlength{\parskip}{0.6em}
\setlength{\parindent}{0pt}

\begin{document}

\begin{center}
{\LARGE \textbf{ECON 300 -- Labour Economics}}\\[0.4em]
{\Large \textbf{Final Exam Review }}\\[0.4em]
Winter 2026
\end{center}

\vspace{0.5em}


% ============================================================
\section*{Part I: Labour Supply and Labour Force Participation}
% ============================================================

\begin{enumerate}

\item A worker has \(T=112\) hours per week to allocate between leisure \(L\) and work \(h\), where
\[
h=T-L.
\]
The worker earns a wage of \(w=\$18\) per hour and has non-labour income of \(\$240\) per week.

\begin{enumerate}[label=(\alph*)]
    \item Write the weekly budget constraint in terms of leisure and income.
    \item What is the intercept on the income axis?
    \item What is the intercept on the leisure axis?
    \item What is the slope of the budget line?
    \item If the wage rises to \(\$24\), show how the budget line changes.
\end{enumerate}

\item A worker's utility-maximizing choice initially is \(L=80\) hours of leisure per week at a wage of \(\$20\) per hour. After a wage increase to \(\$25\), the worker chooses \(L=76\).

\begin{enumerate}[label=(\alph*)]
    \item Compute hours worked before and after the wage increase.
    \item Did labour supply increase or decrease?
    \item Explain, using substitution and income effects, why a wage increase may either increase or decrease hours worked.
\end{enumerate}

\item A person will only enter the labour force if the market wage is at least as high as the reservation wage. Suppose the person values 1 hour of leisure at \(\$14\) and faces fixed costs of working equal to \(\$60\) per day. A full workday is 8 hours.

\begin{enumerate}[label=(\alph*)]
    \item What is the minimum daily earnings needed to just compensate for the value of leisure forgone?
    \item Including fixed costs, what is the minimum daily earnings required to work?
    \item What is the corresponding reservation wage per hour?
    \item If the market wage is \(\$21\), will the person participate?
\end{enumerate}

% ============================================================
\section*{Part II: Income Maintenance, Welfare, and Work Incentives}
% ============================================================

\item Consider an income maintenance program with:
\[
G = 700 \quad \text{and} \quad t = 0.40
\]
where \(G\) is the guaranteed weekly income and \(t\) is the benefit reduction rate. The market wage is \(w=\$25\) per hour.

\begin{enumerate}[label=(\alph*)]
    \item Write disposable income \(C\) as a function of hours worked \(h\) while the worker is on the program.
    \item Compute the effective net wage while receiving benefits.
    \item Find the break-even level of earnings where benefits fall to zero.
    \item How many hours of work does this break-even level correspond to?
    \item Explain how a higher benefit reduction rate affects labour supply incentives.
\end{enumerate}

\item Compare two transfer programs for the same worker:
\begin{itemize}
    \item Program A: demogrant of \(\$500\) per week
    \item Program B: guaranteed income of \(\$500\) with a tax-back rate of \(50\%\)
\end{itemize}
The wage is \(\$20\) per hour.

\begin{enumerate}[label=(\alph*)]
    \item Under Program A, write the budget line equation.
    \item Under Program B, write disposable income while benefits are positive.
    \item Under Program B, what is the effective net wage?
    \item Which program creates a substitution effect? Explain.
    \item Which program is more likely to reduce hours worked? Explain carefully.
\end{enumerate}

% ============================================================
\section*{Part III: Labour Demand in Competitive Markets}
% ============================================================

\item A competitive firm has production function
\[
Q = 40N - N^2
\]
where \(N\) is labour. The output price is \(P=\$3\), and the wage is \(w=\$48\).

\begin{enumerate}[label=(\alph*)]
    \item Compute the marginal product of labour.
    \item Compute the value of the marginal product of labour.
    \item Find the profit-maximizing level of employment.
    \item Compute total output at the optimal employment level.
\end{enumerate}

\item A firm has labour demand elasticity of \(-1.4\). Initially it employs 250 workers at a wage of \(\$30\) per hour. The wage rises by \(10\%\).

\begin{enumerate}[label=(\alph*)]
    \item What is the predicted percentage change in employment?
    \item What is the predicted new level of employment?
    \item Briefly explain what it means for labour demand to be elastic.
\end{enumerate}

\item A firm uses labour \(N\) and capital \(K\). The marginal products are:
\[
MPP_N = 20,\qquad MPP_K = 10.
\]
The wage is \(w=40\), and the rental rate of capital is \(r=15\).

\begin{enumerate}[label=(\alph*)]
    \item Compute \(\frac{MPP_N}{MPP_K}\).
    \item Compute \(\frac{w}{r}\).
    \item Is the cost-minimizing condition satisfied?
    \item Should the firm use relatively more labour or relatively more capital? Explain.
\end{enumerate}

% ============================================================
\section*{Part IV: Monopsony and Wage Determination}
% ============================================================

\item Suppose the labour supply to a monopsonist is
\[
W = 12 + 0.5N
\]
and the marginal revenue product of labour is
\[
MRP_N = 60 - 0.5N.
\]

\begin{enumerate}[label=(\alph*)]
    \item Derive the marginal cost of labour.
    \item Find the monopsony level of employment.
    \item Find the wage paid by the monopsonist.
    \item Find the competitive wage and employment.
    \item Compare the monopsony and competitive outcomes.
\end{enumerate}

\item In a monopsony model, a binding minimum wage of \(\$32\) is introduced. Labour supply remains
\[
W = 12 + 0.5N
\]
and labour demand is
\[
MRP_N = 60 - 0.5N.
\]

\begin{enumerate}[label=(\alph*)]
    \item Find the maximum employment the firm can hire at the minimum wage from the supply curve.
    \item Find the firm's desired employment by setting the relevant marginal cost equal to \(MRP_N\).
    \item Will employment rise or fall relative to the monopsony outcome from the previous question?
    \item Explain why a minimum wage can increase employment under monopsony.
\end{enumerate}

% ============================================================
\section*{Part V: Compensating Wage Differentials}
% ============================================================

\item Two jobs are identical except for workplace risk. Job S is safe and pays \(\$28\) per hour. Job R is risky and pays \(\$34\) per hour.

\begin{enumerate}[label=(\alph*)]
    \item What is the compensating wage differential for risk?
    \item If a worker requires at least \(\$5\) per hour as compensation for risk, which job will the worker choose?
    \item If another worker requires at least \(\$7\) per hour as compensation for risk, which job will that worker choose?
    \item Explain how worker heterogeneity generates sorting across jobs.
\end{enumerate}

\item A worker is willing to accept a dangerous job only if the annual wage premium is at least \(\$1{,}200\). Suppose the job raises annual fatality risk by \(1/20{,}000\).

\begin{enumerate}[label=(\alph*)]
    \item Compute the implied value of a statistical life (VSL).
    \item Explain what this VSL means economically.
\end{enumerate}

% ============================================================
\section*{Part VI: Human Capital, Education, and Training}
% ============================================================

\item Johny is deciding whether to complete a 3-year diploma. The annual costs are:
\begin{itemize}
    \item tuition and books: \(\$2{,}500\)
    \item forgone earnings: \(\$18{,}000\)
\end{itemize}
If he completes the diploma, he expects annual earnings to rise from \(\$30{,}000\) to \(\$38{,}000\) for 20 years after graduation. The discount rate is \(5\%\).

Use the annuity formula:
\[
PV = A\left(\frac{1-(1+r)^{-T}}{r}\right).
\]

\begin{enumerate}[label=(\alph*)]
    \item Compute the total annual cost of education.
    \item Compute the present value of the 3 years of education costs, assuming they are paid at the end of each year.
    \item Compute the annual earnings gain from the diploma.
    \item Compute the present value of the benefits at the time of graduation.
    \item Discount the benefits back to the present.
    \item Should Johny invest in the diploma? Use the NPV rule.
\end{enumerate}

\item A worker receives general training that raises annual earnings by \(\$3{,}000\) in any firm. Another worker receives specific training that raises annual earnings by \(\$4{,}000\) only in the current firm.

\begin{enumerate}[label=(\alph*)]
    \item Which type of training is more likely to be financed by the worker? Why?
    \item Which type is more likely to be financed by the firm? Why?
    \item Explain why the distinction matters for wage-setting during training.
\end{enumerate}

% ============================================================
\section*{Part VII: Discrimination and Oaxaca Decomposition}
% ============================================================

\item The estimated earnings equations are:
\[
\ln W_m = 1.10 + 0.08S + 0.05EXP
\]
\[
\ln W_f = 0.90 + 0.07S + 0.04EXP
\]

Average characteristics are:
\[
\bar S_m = 16,\quad \bar{EXP}_m = 12
\]
\[
\bar S_f = 15,\quad \bar{EXP}_f = 10
\]

\begin{enumerate}[label=(\alph*)]
    \item Compute predicted log earnings for men.
    \item Compute predicted log earnings for women.
    \item Compute the total log wage gap.
    \item Using the male wage structure as the reference, compute the explained portion of the gap:
    \[
    (\bar X_m - \bar X_f)\hat\beta_m
    \]
    \item Compute the unexplained portion.
    \item Interpret the explained and unexplained components.
\end{enumerate}

% ============================================================
\section*{Part VIII: Immigration}
% ============================================================

\item A country admits 360{,}000 immigrants in one year:
\begin{itemize}
    \item Economic: 180{,}000
    \item Family: 108{,}000
    \item Refugees and humanitarian: 72{,}000
\end{itemize}

\begin{enumerate}[label=(\alph*)]
    \item Compute the percentage share of each category.
    \item If economic immigrants rise by 15\% next year while the other categories remain unchanged, compute the new total number of immigrants.
    \item Under the new total, what is the new economic immigrant share?
\end{enumerate}

\item A migrant compares annual net earnings of \(\$48{,}000\) in Country A with \(\$54{,}000\) in Country B. The one-time moving cost is \(\$20{,}000\), the planning horizon is 25 years, and the discount rate is \(4\%\).

\begin{enumerate}[label=(\alph*)]
    \item Compute the annual earnings advantage of Country B.
    \item Compute the present value of that advantage using the annuity formula.
    \item After subtracting moving cost, should the migrant move?
\end{enumerate}

% ============================================================
\section*{Part IX: Unemployment Measurement and Dynamics}
% ============================================================

\item In an economy:
\begin{itemize}
    \item working-age population = 8{,}000{,}000
    \item employed = 5{,}360{,}000
    \item unemployed = 340{,}000
\end{itemize}

\begin{enumerate}[label=(\alph*)]
    \item Compute the labour force.
    \item Compute the unemployment rate.
    \item Compute the labour force participation rate.
    \item Compute the employment rate.
\end{enumerate}

\item Suppose the labour force is 2{,}500{,}000 and the unemployment rate is 8.4\%.

\begin{enumerate}[label=(\alph*)]
    \item How many workers are unemployed?
    \item How many workers are employed?
    \item If 35{,}000 unemployed workers find jobs and 20{,}000 employed workers lose jobs and start searching, what is the new number of unemployed workers?
    \item What is the new unemployment rate?
\end{enumerate}

\item In a region,
\[
UR \approx \text{Incidence} \times \text{Duration}.
\]
Suppose:
\begin{itemize}
    \item Region A: incidence \(= 3\%\), duration \(= 2.5\) months
    \item Region B: incidence \(= 1.5\%\), duration \(= 5\) months
\end{itemize}

\begin{enumerate}[label=(\alph*)]
    \item Compute the unemployment rate in each region.
    \item Which region has higher incidence?
    \item Which region has longer unemployment duration?
    \item Explain why the same unemployment rate can correspond to different policy problems.
\end{enumerate}


\end{enumerate}

\end{document}